\chapter{拉格朗日方程}
\section{系统广义坐标和约束方程}
\( n \) 维广义坐标用 \( n \) 维列向量表示：
\[
\underline{q}(t) = [q_1(t) \ q_2(t) \ \dots \ q_n(t)]^T.
\]
这里要求广义坐标独立，即 \( n = \text{DOFs} \)。
则编号为 \( k \) 的粒子，其全局坐标可表示为
\[
\vec{r}_k(q_1(t), q_2(t), \dots, q_n(t)) =: \vec{r}_k(\underline{q}(t)).
\]
由此我们有系统广义速度：
\[
\begin{aligned}
\frac{d}{dt} \vec{r}_k(\underline{q}(t)) &= \sum_{i=1}^n \frac{\partial \vec{r}_k}{\partial q_i} \frac{d q_i}{dt} \\
&= \sum_{i=1}^n \frac{\partial \vec{r}_k}{\partial q_i} \dot{q}_i \\
&=: \dot{\vec{r}}_k(\underline{q}(t), \dot{\underline{q}}(t)).
\end{aligned}
\]

若 \( \dim(\vec{r}_k) > \text{DOFs} \)，则存在约束方程
\[
0 = C_i(\underline{q}(t), t) \quad (i=1,2,\dots,m).
\]
可写作向量形式
\[
\vec{0} = \vec{C}(\underline{q}(t), t).
\]

下面我们来讨论约束方程对时间的导数：
\[
\begin{aligned}
0 &= C_i(\underline{q}(t), t), \\
0 &= \sum_{j=1}^n \frac{\partial C_i}{\partial q_j} \frac{\partial q_j}{\partial t} + \frac{\partial C_i}{\partial t} \\
&= \sum_{j=1}^n \frac{\partial C_i}{\partial q_j} \dot{q}_j + \frac{\partial C_i}{\partial t} \\
&=: \dot{C}_i(\underline{q}(t), \dot{\underline{q}}(t), t) \quad (i=1,2,\dots,m).
\end{aligned}
\]

若令 \( \vec{C}_{\underline{q}}(\underline{q}(t), t) \)（约束方程的雅各比矩阵）：
\[
\vec{C}_{\underline{q}}(\underline{q}(t), t) = 
\begin{bmatrix}
\frac{\partial C_1}{\partial q_1} & \frac{\partial C_1}{\partial q_2} & \dots & \frac{\partial C_1}{\partial q_n} \\
\frac{\partial C_2}{\partial q_1} & \frac{\partial C_2}{\partial q_2} & \dots & \frac{\partial C_2}{\partial q_n} \\
\vdots & \vdots & \ddots & \vdots \\
\frac{\partial C_m}{\partial q_1} & \frac{\partial C_m}{\partial q_2} & \dots & \frac{\partial C_m}{\partial q_n}
\end{bmatrix} \in \mathbb{R}^{m \times n},
\]
\[
\vec{C}_t(\underline{q}(t), t) = 
\begin{bmatrix}
\frac{\partial C_1}{\partial t} \\
\frac{\partial C_2}{\partial t} \\
\vdots \\
\frac{\partial C_m}{\partial t}
\end{bmatrix} \in \mathbb{R}^{m \times 1},
\]
则有广义速度层面的约束方程：
\[
\vec{0} = \vec{C}_{\underline{q}}(\underline{q}(t), t) \dot{\underline{q}}(t) + \vec{C}_t(\underline{q}(t), t).
\]

\subsection{广义坐标的微小变更}
广义坐标的约束方程：
\[
0 = C_i(\underline{q}(t), t) =: C_i(\underline{q}(t), t).
\]
引入一个广义虚位移 \( \delta \underline{q}(t) \)，
\[
0 = C_i(\underline{q} + \delta \underline{q}, t).
\]
进行一阶 Taylor 展开：
\[
\begin{aligned}
0 &= C_i(\underline{q} + \delta \underline{q}, t) \approx C_i(\underline{q}, t) + \sum_{j=1}^n \frac{\partial C_i(\underline{q}, t)}{\partial q_j} \delta q_j \\
&= \sum_{j=1}^n \frac{\partial C_i(\underline{q}, t)}{\partial q_j} \delta q_j \quad (j=1,2,\dots,n).
\end{aligned}
\]
则
\[
\vec{0} = 
\begin{bmatrix}
\frac{\partial C_1}{\partial q_1} & \frac{\partial C_1}{\partial q_2} & \dots & \frac{\partial C_1}{\partial q_n} \\
\frac{\partial C_2}{\partial q_1} & \frac{\partial C_2}{\partial q_2} & \dots & \frac{\partial C_2}{\partial q_n} \\
\vdots & \vdots & \ddots & \vdots \\
\frac{\partial C_m}{\partial q_1} & \frac{\partial C_m}{\partial q_2} & \dots & \frac{\partial C_m}{\partial q_n}
\end{bmatrix}
\begin{bmatrix}
\delta q_1 \\
\delta q_2 \\
\vdots \\
\delta q_n
\end{bmatrix}
=: \vec{C}_{\underline{q}}(\underline{q}, t) \delta \underline{q} = \frac{\partial \vec{C}(\underline{q}, t)}{\partial \underline{q}} \delta \underline{q}.
\]

\subsection{质点的虚位移}
\[
\vec{r}_k = \vec{r}_k(\underline{q}(t)), \quad \vec{r}_k(\underline{q} + \delta \underline{q}),
\]
\[
\delta \vec{r}_k = \vec{r}_k(\underline{q} + \delta \underline{q}) - \vec{r}_k(\underline{q}).
\]
由一阶 Taylor 展开：
\[
\delta \vec{r}_k \approx \sum_{j=1}^n \frac{\partial \vec{r}_k}{\partial q_j} \delta q_j,
\]
即
\[
\delta \vec{r}_k = \sum_{j=1}^n \frac{\partial \vec{r}_k}{\partial q_j} \delta q_j.
\]

\section{虚功原理与拉格朗日方程}
\subsection{质点的运动微分方程}
Newton 第二运动定律：
\[
\vec{F} = \frac{d}{dt}(m \vec{v}) = \frac{d}{dt}\left(m \frac{d \vec{r}}{dt}\right) = \frac{dm}{dt} \frac{d \vec{r}}{dt} + m \frac{d^2 \vec{r}}{dt^2}.
\]
一个具有恒定质量的物体（不考虑相对论效应），则 \( \frac{dm}{dt} = 0 \)，
\[
\vec{F} = m \frac{d^2 \vec{r}}{dt^2} =: m \ddot{\vec{r}} =: m \vec{a}.
\]
若将 \( (-m \ddot{\vec{r}}) \) 视作惯性力（Inertial force），则有 D'Alembert 原理（达朗贝尔原理）：
\[
-m \ddot{\vec{r}} + \vec{F} = \vec{0}.
\]
对于一个由 \( k \) 个部分组成的多体系统：
\[
-m_k \ddot{\vec{r}}_k + \vec{F}_{{\rm EX},k} + \vec{F}_{{\rm IN},k} = \vec{0}.
\]

\subsection{质点与质点系的虚功方程}
对于质点系中的第 \( k \) 个质点，有动力学方程
\[
-m_k \ddot{\vec{r}}_k + \vec{F}_{{\rm EX},k} + \vec{F}_{{\rm IN},k} = \vec{0}.
\]
由此有虚功方程
\[
-m_k \delta \vec{r}_k^T \ddot{\vec{r}}_k + \delta \vec{r}_k^T \vec{F}_{{\rm EX},k} + \delta \vec{r}_k^T \vec{F}_{{\rm IN},k} = \vec{0}.
\]
将质点系中所有的方程相加：
\[
-\sum_{k=1}^N m_k \delta \vec{r}_k^T \ddot{\vec{r}}_k + \sum_{k=1}^N \delta \vec{r}_k^T \vec{F}_{{\rm EX},k} + \sum_{k=1}^N \delta \vec{r}_k^T \vec{F}_{{\rm IN},k} = \vec{0}.
\]
可以看到，虚功方程由三部分组成：
\begin{enumerate}
\item 惯性力的虚功：\( -\sum_{k=1}^N m_k \delta \vec{r}_k^T \ddot{\vec{r}}_k \)
\item 系统外力的虚功：\( \sum_{k=1}^N \delta \vec{r}_k^T \vec{F}_{{\rm EX},k} \)
\item 系统内力的虚功：\( \sum_{k=1}^N \delta \vec{r}_k^T \vec{F}_{{\rm IN},k} \)
\end{enumerate}

下面逐一剖析：

\subsubsection{系统惯性力的虚功}
\[
\begin{aligned}
-\sum_{k=1}^N m_k \delta \vec{r}_k^T \ddot{\vec{r}}_k &= -\sum_{k=1}^N m_k \left( \sum_{j=1}^n \frac{\partial \vec{r}_k}{\partial q_j} \delta q_j \right)^T \frac{d}{dt}\left( \dot{\vec{r}}_k \right) \\
&= -\sum_{k=1}^N m_k \sum_{j=1}^n \delta q_j \frac{\partial \vec{r}_k^T}{\partial q_j} \frac{d}{dt}\left( \dot{\vec{r}}_k \right) \\
&= -\sum_{j=1}^n \sum_{k=1}^N m_k \delta q_j \frac{\partial \vec{r}_k^T}{\partial q_j} \frac{d}{dt}\left( \dot{\vec{r}}_k \right) \\
&= -\sum_{j=1}^n \delta q_j \sum_{k=1}^N m_k \left[ \frac{d}{dt}\left( \frac{\partial \vec{r}_k^T}{\partial q_j} \dot{\vec{r}}_k \right) - \frac{d}{dt}\left( \frac{\partial \vec{r}_k^T}{\partial q_j} \right) \dot{\vec{r}}_k \right].
\end{aligned}
\]
注意到
\[
\dot{\vec{r}}_k = \frac{d \vec{r}_k}{dt} = \sum_{j=1}^n \frac{\partial \vec{r}_k}{\partial q_j} \dot{q}_j + \frac{\partial \vec{r}_k}{\partial t},
\]
\[
\frac{\partial \dot{\vec{r}}_k}{\partial \dot{q}_j} = \frac{\partial}{\partial \dot{q}_j}\left( \sum_{j=1}^n \frac{\partial \vec{r}_k}{\partial q_j} \dot{q}_j \right) = \frac{\partial \vec{r}_k}{\partial q_j},
\]
即
\[
\frac{\partial \dot{\vec{r}}_k}{\partial \dot{q}_j} = \frac{\partial \vec{r}_k}{\partial q_j}.
\]
又
\[
\frac{d}{dt}\left( \frac{\partial \vec{r}_k}{\partial q_j} \right) = \frac{\partial \dot{\vec{r}}_k}{\partial q_j}.
\]
于是
\[
\begin{aligned}
\sum_{k=1}^N m_k \frac{\partial \vec{r}_k^T}{\partial q_j} \dot{\vec{r}}_k &= \sum_{k=1}^N m_k \frac{\partial \dot{\vec{r}}_k^T}{\partial \dot{q}_j} \dot{\vec{r}}_k \\
&= \sum_{k=1}^N m_k \frac{\partial}{\partial \dot{q}_j} \left( \frac{1}{2} \dot{\vec{r}}_k^T \dot{\vec{r}}_k \right) \\
&= \frac{\partial}{\partial \dot{q}_j} \left( \sum_{k=1}^N \frac{1}{2} m_k \dot{\vec{r}}_k^T \dot{\vec{r}}_k \right) = \frac{\partial T}{\partial \dot{q}_j}.
\end{aligned}
\]
因此
\[
\begin{aligned}
-\sum_{k=1}^N m_k \delta \vec{r}_k^T \ddot{\vec{r}}_k &= -\sum_{j=1}^n \delta q_j \left[ \frac{d}{dt}\left( \frac{\partial T}{\partial \dot{q}_j} \right) - \frac{\partial T}{\partial q_j} \right] \\
&= \sum_{j=1}^n \delta q_j \left[ -\frac{d}{dt}\left( \frac{\partial T}{\partial \dot{q}_j} \right) + \frac{\partial T}{\partial q_j} \right].
\end{aligned}
\]

\subsubsection{系统内力与外力的虚功}
\[
\begin{aligned}
\sum_{k=1}^N \delta \vec{r}_k^T \vec{F}_{{\rm EX},k} + \sum_{k=1}^N \delta \vec{r}_k^T \vec{F}_{{\rm IN},k}
&= \sum_{k=1}^N \delta \vec{r}_k^T \vec{F}_{{\rm EX},k} + \sum_{k=1}^N \delta \vec{r}_k^T \vec{F}_{{\rm NonIdeal},k} \\
&= \sum_{k=1}^N \delta \vec{r}_k^T \left( \vec{F}_{{\rm EX},k} + \vec{F}_{{\rm NonIdeal},k} \right) \\
&= \sum_{k=1}^N \sum_{j=1}^n \delta q_j \frac{\partial \vec{r}_k^T}{\partial q_j} \left( \vec{F}_{{\rm EX},k} + \vec{F}_{{\rm NonIdeal},k} \right) \\
&= \sum_{j=1}^n \delta q_j \sum_{k=1}^N \frac{\partial \vec{r}_k^T}{\partial q_j} \left( \vec{F}_{{\rm EX},k} + \vec{F}_{{\rm NonIdeal},k} \right) \\
&=: \sum_{j=1}^n \delta q_j Q_j.
\end{aligned}
\]

于是，虚功方程被简化为：
\[
0 = \sum_{j=1}^n \delta q_j \left[ -\frac{d}{dt}\left( \frac{\partial T}{\partial \dot{q}_j} \right) + \frac{\partial T}{\partial q_j} + Q_j \right].
\]
写成矩阵形式：
\[
0 = \delta \underline{q}^T \left[ -\frac{d}{dt}\left( \frac{\partial T}{\partial \dot{\underline{q}}} \right)^T + \left( \frac{\partial T}{\partial \underline{q}} \right)^T + \vec{Q} \right],
\]
其中 \( \delta \underline{q} \) 满足 \( \underline{0} = \frac{\partial \vec{C}}{\partial \underline{q}} \delta \underline{q} \)。
由广义变分原理，\( \exists \underline{\lambda} \in \mathbb{R}^m \)，s.t.
\[
-\frac{d}{dt}\left( \frac{\partial T}{\partial \dot{\underline{q}}} \right)^T + \left( \frac{\partial T}{\partial \underline{q}} \right)^T + \vec{Q} - \left( \frac{\partial \vec{C}}{\partial \underline{q}} \right)^T \underline{\lambda} = \underline{0}.
\]

\section{系统动能与广义质量矩阵}
\subsection{系统动能 (The system kinetic energy)}
\begin{definition}[系统动能]
\[
T = \sum_{k=1}^N \frac{1}{2} m_k \dot{\vec{r}}_k^T \dot{\vec{r}}_k.
\]
这里，
\[
\dot{\vec{r}}_k = \sum_{j=1}^n \frac{\partial \vec{r}_k}{\partial q_j} \frac{d q_j}{dt} = \sum_{j=1}^n \frac{\partial \vec{r}_k}{\partial q_j} \dot{q}_j =: \vec{r}_{k,\underline{q}} \dot{\underline{q}}.
\]
则
\[
\begin{aligned}
T &= \sum_{k=1}^N \frac{1}{2} m_k \left( \vec{r}_{k,\underline{q}} \dot{\underline{q}} \right)^T \left( \vec{r}_{k,\underline{q}} \dot{\underline{q}} \right) \\
&= \sum_{k=1}^N \frac{1}{2} m_k \dot{\underline{q}}^T \vec{r}_{k,\underline{q}}^T \vec{r}_{k,\underline{q}} \dot{\underline{q}} \\
&= \frac{1}{2} \dot{\underline{q}}^T \left( \sum_{k=1}^N m_k \vec{r}_{k,\underline{q}}^T \vec{r}_{k,\underline{q}} \right) \dot{\underline{q}}.
\end{aligned}
\]
\end{definition}

\subsection{系统广义质量矩阵 (The generalized inertial matrix)}
\begin{definition}[广义质量矩阵]
\[
\vec{M}(\underline{q}) = \sum_{k=1}^N m_k \vec{r}_{k,\underline{q}}^T \vec{r}_{k,\underline{q}}.
\]
则系统动能
\[
T = \frac{1}{2} \dot{\underline{q}}^T \vec{M}(\underline{q}) \dot{\underline{q}}.
\]
\end{definition}

\begin{proposition}[广义质量矩阵的性质]
\begin{enumerate}
\item \textbf{对称性 (Symmetry)}：\( \vec{M} = \vec{M}^T \)。
\begin{proof}
\[
M_{jk} = \sum_{i=1}^N m_i \frac{\partial \vec{r}_i}{\partial q_j} \cdot \frac{\partial \vec{r}_i}{\partial q_k}, \quad M_{kj} = \sum_{i=1}^N m_i \frac{\partial \vec{r}_i}{\partial q_k} \cdot \frac{\partial \vec{r}_i}{\partial q_j}.
\]
显然 \( M_{jk} = M_{kj} \)，故 \( \vec{M} = \vec{M}^T \)。
\end{proof}

\item \textbf{正定性 (Positive Definiteness)}：\( \vec{M} \) 始终正定或半正定。
\begin{proof}
\( T \geq 0 \implies 2T \geq 0 \)，即
\[
\dot{\underline{q}}^T \vec{M}(\underline{q}) \dot{\underline{q}} \geq 0.
\]
故 \( \vec{M}(\underline{q}) \) 正定或半正定。
\end{proof}
\end{enumerate}
\end{proposition}

\begin{remark}
存在奇异位姿，即 \( |\vec{M}| = 0 \) 时，\( \vec{M}(\underline{q}) \) 半正定。
\end{remark}